%%%%%%%%%%%%%%%%%%%%%%%%%%%%%%%%%%%%%%%%%%%%%%%%%%%%%%%%%%%%%%%%%%%%%%%%%%%%%
%  enslyonreview-gabarit.tex — blank template for the M2 Biosciences
%  bibliographic review (UE "Analyse bibliographique", Current Opinion style)
%
%  New to LaTeX? Everything between \begin{document} and \end{document} is
%  what actually gets printed; everything above \begin{document} (the
%  "preamble") only configures the document and never appears in the PDF.
%  A command looks like \commandname{argument} — curly braces {} wrap the
%  text you give it; a "%" starts a comment (like this line) and is ignored
%  by LaTeX, so feel free to leave these notes in place while you write.
%
%  → Fill in the fields below, then write your sections.
%  → For a fully worked example, open enslyonreview-exemple.tex (same
%    structure, real content: figures, a table, 50 cited references).
%  → Full guide, with the rendered result of every command: enslyonreview.pdf
%
%  On Overleaf: this file is detected automatically as the root document,
%  and biber (the bibliography processor) runs automatically on every
%  compile.
%  Locally: pdflatex enslyonreview-gabarit → biber enslyonreview-gabarit
%           → pdflatex enslyonreview-gabarit → pdflatex enslyonreview-gabarit
%  (Yes, four steps — that repetition is normal for a document with a
%  bibliography; see enslyonreview.pdf, section "Compilation", for why.)
%
%  Project files:
%    enslyonreview-gabarit.tex ← this file (fill in and compile)
%    enslyonreview-exemple.tex ← fully worked example (for reference)
%    enslyonreview.cls         ← the class (same folder, or on TEXMFHOME)
%    enslyonreview.dbx         ← extra bibliography fields (pmid/pmcid/
%                                 interest) — keep next to the .cls,
%                                 never edit it
%    references.bib            ← your reference database (a Zotero export)
%    figures/                  ← your images (pdf, png, jpg...)
%%%%%%%%%%%%%%%%%%%%%%%%%%%%%%%%%%%%%%%%%%%%%%%%%%%%%%%%%%%%%%%%%%%%%%%%%%%%%

% ═══════════════════════════════════════════════════════════════════════════
%  THE "CURRENT OPINION" STYLE — what this class reproduces
% ═══════════════════════════════════════════════════════════════════════════
%  Conventions observed in the review_M2.ott template and in an already
%  finished, approved review (not rules LaTeX forces on you, but the format
%  this UE expects):
%    ✓ Section headings with NO number ("Introduction", not "1. Introduction")
%    ✓ Full-width title block: title, author, affiliation/date, keywords,
%      abstract — printed before the text switches to two columns
%    ✓ Abstract as a single paragraph, with NO bibliographic reference
%    ✓ Citations as bracketed numbers [1], compressed into ranges like
%      [8-10] for consecutive references
%    ✓ References flagged "of special interest" (•) / "of outstanding
%      interest" (••), each with a one-sentence note on why it matters
%    ✓ Figures/tables spanning the full width (figure*/table*) or a single
%      column (figure/table), plus "Box N." call-out panels for a side point
%  Every command below, with what it actually looks like once compiled, is
%  explained in enslyonreview.pdf.
% ═══════════════════════════════════════════════════════════════════════════

% ═══════════════════════════════════════════════════════════════════════════
%  UE "ANALYSE BIBLIOGRAPHIQUE" — OFFICIAL GUIDELINES  (keep these comments
%  during writing; source: F. Roudier's "Guidelines for the bibliographic
%  analysis course", 2025–2026 cohort — confirm with him that the numbers
%  below still hold for your own year before submitting)
% ═══════════════════════════════════════════════════════════════════════════
%  SUBJECT
%    ✓ Connected to your internship's research topic, but with a LARGER
%      scope (not redundant with your internship report's introduction)
%    ✓ A field broad enough for a significant bibliographic analysis
%    ✓ Recent scientific advances; not covered by a recently published review
%    ✓ Compatible with the publication standards of review articles such as
%      Current Opinion (which is exactly what this class reproduces)
%    ✓ Proposed with your internship supervisor, then validated by F. Roudier
%
%  LENGTH & DISPLAY ITEMS
%    30-50 references · ~25 000 characters · 3-5 figures/tables
%    → "Specific guidelines to format the review, the artwork and the
%      references" are announced as being communicated separately/later
%      each year — review_M2.ott and this class follow the format of an
%      already-approved review (see enslyonreview-exemple.tex) in the
%      absence of a more detailed written spec
%
%  SUBMISSION
%    PDF, submitted on the Portal — check the CURRENT year's exact deadline
%    with F. Roudier
% ═══════════════════════════════════════════════════════════════════════════

% ── COLOUR THEME ─────────────────────────────────────────────────────────
% This is the very first command a LaTeX document needs: \documentclass
% tells LaTeX which class (i.e. which set of formatting rules) to use — here,
% enslyonreview. The optional [name] in square brackets is a "class option";
% for this class it picks a colour theme.
%
% 23 themes are available; "orange" (the default — used when you pass no
% option at all, as below) is the official rust/peach colour scheme of the
% .ott template — keep it for your final version. "orangeens" is the
% official ENS de Lyon institutional orange instead (different from the
% .ott's own orange — same colour as enslyonstage's default theme). The
% other 21 are just for your own experiments:
%   orangeens · violet · blue · red · green · teal · gray · navy · azure ·
%   indigo · midnight · slate · purple · pink · crimson · wine · amber ·
%   gold · brown · lime · forest · cyan
\documentclass{enslyonreview}
% \documentclass[blue]{enslyonreview}       % ← uncomment to try another theme
% \documentclass[palatino]{enslyonreview}   % ← Palatino instead of Times (newtx)
% \documentclass[xcharter]{enslyonreview}   % ← Charter instead of Times (newtx)
% \documentclass[stix]{enslyonreview}       % ← STIX Two instead of Times (newtx)
% A theme and a font option can be combined: \documentclass[blue, stix]{enslyonreview}

% ── BIBLIOGRAPHY ────────────────────────────────────────────────────────────
% This tells LaTeX which file holds your references. It uses biblatex with
% the biber backend — compile with biber (see the compile sequence at the
% top of this file), NOT with the older "bibtex" command.
\addbibresource{references.bib}

% ── "OF SPECIAL / OUTSTANDING INTEREST" REFERENCES (optional) ──────────────
% Declared HERE, in this .tex file — NOT inside the .bib: a normal Zotero
% export/re-export does not carry this annotation along, and re-exporting
% later would silently wipe out anything you had added by hand to the .bib.
%   \refspecialinterest{citationKey}{Why this reference matters.}
%   \refoutstandinginterest{anotherKey}{Why this one matters even more.}
% The key must match EXACTLY (including capitalisation) the key of that
% entry in your .bib file. These two commands can go anywhere before
% \printbibliography — right here, under \addbibresource, is the easiest
% place to keep track of them.


%%%%%%%%%%%%%%%%%%%%%%%%%%%%%%%%%%%%%%%%%%%%%%%%%%%%%%%%%%%%%%%%%%%%%%%%%%%%%
%  METADATA — fill in every field below before compiling
%  These feed the full-width title block (page 1), the running header on
%  every page, and the footer.
%%%%%%%%%%%%%%%%%%%%%%%%%%%%%%%%%%%%%%%%%%%%%%%%%%%%%%%%%%%%%%%%%%%%%%%%%%%%%

% ── TITLE ───────────────────────────────────────────────────────────────
% The full title, printed at the top of page 1.
\reviewtitle{Your full title here}

% The short title, printed in the running header on EVERY page — a few
% words are enough; the full title above wouldn't fit on one line there.
\reviewshorttitle{Short title}

% ── AUTHOR ──────────────────────────────────────────────────────────────
% \textsc{} renders text in SMALL CAPITALS — used here on the surname
% (same convention as enslyonstage). It works identically with the default
% font and with any of the [palatino]/[xcharter]/[stix] options above.
\reviewauthor{First name \textsc{Surname}}
% Short form (usually just the surname), shown in the running header next
% to the short title.
\reviewauthorshort{Surname}

% ── AFFILIATION AND DATE ────────────────────────────────────────────────
% These two are printed together, in the peach band under the title; the
% date alone is repeated in the footer of every page.
\reviewaffiliation{Master Biosciences, D\'epartement de Biologie, \'Ecole Normale Sup\'erieure de Lyon.}
\reviewdate{YYYY-MM-DD}

% ── KEYWORDS ────────────────────────────────────────────────────────────
% A handful of keywords, separated by commas — 4 to 6 is a typical length
% in the supplied example, but the class doesn't enforce any count.
\reviewkeywords{keyword1, keyword2, keyword3, keyword4}

% ── ABSTRACT ────────────────────────────────────────────────────────────
% A single paragraph, with NO bibliographic reference — a Current Opinion
% convention: none of the reference reviews cite a source in their abstract.
\reviewabstract{Your abstract here, as a single paragraph, with no
  bibliographic reference.}


%%%%%%%%%%%%%%%%%%%%%%%%%%%%%%%%%%%%%%%%%%%%%%%%%%%%%%%%%%%%%%%%%%%%%%%%%%%%%
\begin{document}
%%%%%%%%%%%%%%%%%%%%%%%%%%%%%%%%%%%%%%%%%%%%%%%%%%%%%%%%%%%%%%%%%%%%%%%%%%%%%
% Everything from here down is what actually appears in the PDF.

% ── TITLE BLOCK — builds the full-width title/author/affiliation/abstract
% block above, THEN automatically switches to two columns for everything
% that follows. This must stay the very first command after \begin{document}.
\maketitle

%%%%%%%%%%%%%%%%%%%%%%%%%%%%%%%%%%%%%%%%%%%%%%%%%%%%%%%%%%%
\section{Introduction}
%%%%%%%%%%%%%%%%%%%%%%%%%%%%%%%%%%%%%%%%%%%%%%%%%%%%%%%%%%%
% You don't strictly need a \label for the Introduction (it's rarely
% referenced from elsewhere in the document) — but do add one to any
% section you might want to point back to later, e.g. \label{sec:mechanisms},
% then refer to it with \cref{sec:mechanisms} (or \crefcolor{} for a bold,
% theme-coloured version) instead of typing "see section 2" by hand — that
% way the number stays correct even if you reorder your sections later.
%
% Cite a reference with \verb|\cite{citationKey}| — the numbers are
% assigned, and compressed into ranges ([8-10]), automatically, in the
% order each reference is first cited in the text. The key is the one used
% for that entry in your .bib file (e.g. a Zotero export).

Your introduction text, with citations to your sources\ldots


%%%%%%%%%%%%%%%%%%%%%%%%%%%%%%%%%%%%%%%%%%%%%%%%%%%%%%%%%%%
\section{Your first part}
\label{sec:part1}
%%%%%%%%%%%%%%%%%%%%%%%%%%%%%%%%%%%%%%%%%%%%%%%%%%%%%%%%%%%
% Each major section typically covers one distinct mechanism or theme (see
% enslyonreview-exemple.tex: "Nucleotide or nucleic acid degradation",
% "Transcription inhibition"...). \subsection{} works exactly like
% \section{} inside it, still with no number shown.


%%%%%%%%%%%%%%%%%%%%%%%%%%%%%%%%%%%%%%%%%%%%%%%%%%%%%%%%%%%
\section{Your second part}
%%%%%%%%%%%%%%%%%%%%%%%%%%%%%%%%%%%%%%%%%%%%%%%%%%%%%%%%%%%


% ── FULL-WIDTH FIGURE (spans both columns) ────────────────────────────────
% figure* spans both columns, plain figure stays within one — standard
% LaTeX behaviour in two-column mode, no special macro to learn. A table
% follows exactly the same logic with table/table*. See enslyonreview.pdf
% (section "Figures and tables") for the other variants and a real
% rendered example.
%\begin{figure*}[!t]
%  \centering
%  \includegraphics[width=0.9\textwidth]{figures/myimage}
%  \caption{\textbf{Short title in bold.} The rest of the caption, in
%    normal text — self-sufficient: a reader should understand it without
%    having to re-read the body text.}
%  \label{fig:myimage}
%\end{figure*}


%%%%%%%%%%%%%%%%%%%%%%%%%%%%%%%%%%%%%%%%%%%%%%%%%%%%%%%%%%%
\section{Concluding remarks}
%%%%%%%%%%%%%%%%%%%%%%%%%%%%%%%%%%%%%%%%%%%%%%%%%%%%%%%%%%%
% The usual closing section in the Current Opinion format — restates the
% key takeaways and opens onto future perspectives, without introducing any
% new citation-backed claim.


%%%%%%%%%%%%%%%%%%%%%%%%%%%%%%%%%%%%%%%%%%%%%%%%%%%%%%%%%%%
\section*{Acknowledgements}
%%%%%%%%%%%%%%%%%%%%%%%%%%%%%%%%%%%%%%%%%%%%%%%%%%%%%%%%%%%
% The star in \section*{} means: an unnumbered section that also stays out
% of any \tableofcontents you might add later — appropriate for an
% "outside the outline" section like this one.


%--- Bibliography ----------------------------------------------------------
% \balance evens out the height of the two columns on the very last page
% (so one column doesn't end noticeably higher than the other) — call it
% right after \printbibliography, as in enslyonreview-exemple.tex, never
% earlier (doing so could disturb the layout of earlier pages).
\printbibliography[title=References]
% Or \printbibliography[title=References and recommended reading] to keep
% the heading suggested by the .ott — that's the default title anyway if
% you omit the title= option entirely.

%\balance

\end{document}
